% Ensure included PDFs (banners) with newer versions embed without warnings
\pdfminorversion=7
\documentclass[aspectratio=169]{beamer}

\makeatletter
% Make LaTeX find theme .sty files in Template
\def\input@path{{Template/}}
\makeatother
\usepackage{graphicx}
% Make graphics (including banner PDFs) resolvable without TEXINPUTS
\graphicspath{{Template/}{Template/banners/}{../Common_Images/}}

%================================================================%
% Theme and Package Setup
%================================================================%
\usetheme{ucl}
\setbeamercolor{banner}{bg=darkpurple}

% Navigation and Footer
\setbeamertemplate{navigation symbols}{\vspace{-2ex}}
\setbeamertemplate{footline}[author title date]
\setbeamertemplate{slide counter}[framenumber/totalframenumber]

\usepackage[utf8]{inputenc}
\usepackage[british]{babel}
\usepackage{amsmath, amssymb, amsthm}
\usepackage{tikz}
\usepackage{pgfplots}
\pgfplotsset{compat=1.17}
\usetikzlibrary{calc, positioning, arrows.meta, shapes.geometric, trees, backgrounds, shapes.misc, graphs, quotes, shadows, fit, decorations.pathreplacing} 
\usefonttheme{professionalfonts}
\usepackage{eulervm}
\usepackage{listings}
\usepackage{xcolor}
\usepackage{algorithm}
\usepackage{algpseudocode}
\usepackage[square,authoryear]{natbib}

% Define custom colours
\makeatletter
\@ifundefined{color@stone}{%
    \definecolor{stone}{gray}{0.95}%
}{}
\makeatother
\definecolor{darkgreen}{rgb}{0.0, 0.5, 0.0}
\definecolor{uclblue}{cmyk}{1,0.24,0,0.64}

\setbeamercovered{transparent}

% Define theoremblock
\makeatletter
\newenvironment<>{theoremblock}[1]{%
    \begin{block}#2{#1}%
}{\end{block}}
\makeatother

% Section divider slides
\AtBeginSection[]{
  \begin{frame}
    \frametitle{\textbf{\Large\insertsectionhead}}
    \begin{center}
        \huge\textbf{\insertsectionhead}
    \end{center}
  \end{frame}
}

% Operator definitions
\DeclareMathOperator*{\argmin}{arg\,min}

\title{Foundations of Deep Learning}
\subtitle{From Perceptrons to Continuous Dynamics and Implicit Models}
\author[Martin Benning (University College London)]{Martin Benning}
\date[CCMI]{CCMI -- Deep Learning in Theory and Practice \\[1cm] Monday, 18 May 2026}
\institute[]{University College London}

\begin{document}

% --- Title Frame ---
\begin{frame}
  \titlepage
\end{frame}

% --- Course Overview ---
\begin{frame}{Course Overview}
    \begin{block}{Course Description}
        This course provides a comprehensive overview of deep learning theory and its application to inverse problems, bridging the gap between classical mathematical methods and modern data-driven architectures.
    \end{block}
    
    \vspace{2em}
    \begin{center}
        \Large \textbf{Course Materials:} \\
        \vspace{0.5em}
        \normalsize \url{https://github.com/Deep-Learning-Theory-Practice-2026/Course-Materials-Public}
    \end{center}
\end{frame}

% --- Weekly Schedule ---
\begin{frame}{Weekly Schedule}
    \begin{itemize}
        \item \textbf{Monday:} Foundations of Deep Learning
        \begin{itemize}
            \item Architectures, training dynamics, and backpropagation.
            \item Pitfalls of end-to-end learning for inverse problems.
        \end{itemize}
        \item \textbf{Tuesday:} Unrolling, Bilevel Optimisation \& Physics-Informed Learning
        \begin{itemize}
            \item Null-space networks, variational methods, and optimisation.
            \item Bilevel learning via implicit differentiation and algorithm unrolling (LISTA, LPD).
        \end{itemize}
        \item \textbf{Wednesday:} Modular Deep Learning, PnP \& Generative Models
        \begin{itemize}
            \item Plug-and-Play (PnP) priors, RED, and score-based diffusion models.
        \end{itemize}
        \item \textbf{Thursday:} Continuous \& Implicit Deep Learning
        \begin{itemize}
            \item Neural ODEs, optimal control, and Deep Equilibrium Models (DEQs).
        \end{itemize}
        \item \textbf{Friday:} Coding Challenge Presentations
    \end{itemize}
\end{frame}

% --- Overview Frame ---
\begin{frame}{Lecture Overview}
  \begin{columns}[T]
    \column{0.5\textwidth}
    \begin{block}{Objectives}
        \begin{itemize}
            \item Foundational Architectures (Perceptrons, MLPs, CNNs).
            \item Training Dynamics, Optimisation, and Backpropagation.
            \item Advanced Architectures (ResNets, UNets).
            \item Generalised Backpropagation for Implicit \& Continuous Models.
            \item Pitfalls of End-to-End Learning for Inverse Problems.
        \end{itemize}
    \end{block}
    \column{0.5\textwidth}
    \tableofcontents
  \end{columns}
\end{frame}

%================================================================%
\section{The Deep Learning Revolution \& Motivation}
%================================================================%

\begin{frame}{The Rise of Deep Learning}
    \begin{columns}
        \column{0.5\textwidth}
        \textbf{Deep Learning is everywhere:}
        \begin{itemize}
            \item Image Classification (ImageNet).
            \item Speech Recognition (Siri, Alexa).
            \item Natural Language Processing (Transformers, GPT).
            \item Strategic Games (AlphaGo).
        \end{itemize}
        
        \column{0.5\textwidth}
        \centering
        \begin{tikzpicture}[scale=0.6]
            \begin{axis}[
                title={ImageNet Top-5 Error Rate},
                xlabel={Year},
                ylabel={Error Rate (\%)},
                xmin=2010, xmax=2017,
                ymin=0, ymax=30,
                xtick={2010, 2011, 2012, 2013, 2014, 2015, 2016},
                xticklabels={2010, 2011, 2012, 2013, 2014, 2015, 2016},
                ymajorgrids=true,
                grid style=dashed,
            ]
            \addplot[
                color=red,
                mark=square,
                thick
                ]
                coordinates {
                (2010,28.2)(2011,25.8)(2012,16.4)(2013,11.7)(2014,6.7)(2015,3.6)(2016,3.0)
                };
            \node at (axis cs:2012,18) [anchor=south west] {AlexNet};
            \node at (axis cs:2015,5) [anchor=south west] {ResNet};
            \draw[dashed, blue] (axis cs:2010,5.1) -- (axis cs:2017,5.1) node[midway, below] {Human Performance};
            \end{axis}
        \end{tikzpicture}
    \end{columns}
\end{frame}

%================================================================%
\section{Deep Learning Primitives: Architectures \& Activations}
%================================================================%

\begin{frame}{The Artificial Neuron}
    The fundamental building block of deep learning is the \textbf{Artificial Neuron} (with its special case the \textbf{Perceptron} \citep{rosenblatt1957,rosenblatt1958,rosenblatt1962}).\vspace{0.1cm}
    
    \begin{columns}
        \column{0.4\textwidth}
        Mathematical formulation:\vspace{-0.1cm}
        $$ y = \sigma\left( \sum_{i=1}^n w_i x_i + b \right) $$
        \begin{itemize}
            \item $x$: Inputs
            \item $w$: Weights (Synaptic strength)
            \item $b$: Bias (Threshold)
            \item $\sigma$: Activation Function (Firing rate)
        \end{itemize}
        
        \column{0.6\textwidth}
        \centering
        \begin{tikzpicture}[
            init/.style={
              circle,
              draw=black,
              fill=stone,
              minimum size=0.5cm
            },
            nnode/.style={
              circle,
              draw=black,
              minimum size=1cm,
              fill=uclblue!20
            }
        ]
        
        % Inputs
        \node[init] (x1) at (0,2) {$x_1$};
        \node[init] (x2) at (0,1) {$x_2$};
        \node[init] (xn) at (0,-1) {$x_n$};
        \node at (0,0.2) {$\vdots$};
        
        % Weights labels
        \node (w1) at (1.5,2.2) {$w_1$};
        \node (w2) at (1.5,1.2) {$w_2$};
        \node (wn) at (1.5,-0.8) {$w_n$};
        
        % Neuron
        \node[nnode] (sigma) at (3,0.5) {$\sum$};
        \node[nnode] (act) at (5,0.5) {$\sigma(\cdot)$};
        
        % Output
        \node (out) at (7,0.5) {$y$};
        
        % Connections
        \draw[->] (x1) -- (sigma);
        \draw[->] (x2) -- (sigma);
        \draw[->] (xn) -- (sigma);
        \draw[->] (sigma) -- (act) node[midway, above] {$z$};
        \draw[->] (act) -- (out);
        
        \node[above] at (sigma.north) {Aggregation};
        \node[above] at (act.north) {Activation};
        \end{tikzpicture}
    \end{columns}
\end{frame}

% \begin{frame}{Solving decision problems}
%     \begin{itemize}
%         \item Explain how to model decision problems
%         \item Include concert example
%         \item Model with Heaviside function 
%     \end{itemize}
% \end{frame}

\begin{frame}{Activation Functions}
    The power of neural networks comes from \textbf{non-linearity}. Without $\sigma$, a deep network is just a single linear matrix multiplication.
    
    \begin{columns}[T]
        \column{0.33\textwidth}
        \centering
        \textbf{Sigmoid}
        \vspace{-0.35cm}
        
        $$ \sigma(z) = \frac{1}{1+e^{-z}} $$
        \vspace{-0.15cm}
        
        \begin{tikzpicture}[scale=0.55]
            \begin{axis}[width=5.5cm, height=5cm, axis lines=middle, ymin=-0.1, ymax=1.1]
                \addplot[blue, thick, domain=-5:5, samples=100] {1/(1+exp(-x))};
            \end{axis}
        \end{tikzpicture}
        
        \vspace{0.05cm}
        
        \footnotesize{Historically used, but suffers from vanishing gradients.}

        \column{0.33\textwidth}
        \centering
        \textbf{ReLU}
        \vspace{-0.35cm}
        
        $$ \sigma(z) = \max(0, z) $$
        \vspace{-0.15cm}
        
        \begin{tikzpicture}[scale=0.55]
            \begin{axis}[width=5.5cm, height=5cm, axis lines=middle, ymin=-1, ymax=5]
                \addplot[red, thick, domain=-5:5, samples=100] {max(0,x)};
            \end{axis}
        \end{tikzpicture}
        
        \vspace{0.05cm}
        
        \footnotesize{Standard in modern DL. Sparse activation, no vanishing gradient for $z>0$.}

        \column{0.33\textwidth}
        \centering
        \textbf{Leaky ReLU}
        \vspace{-0.35cm}
        
        $$ \max(0.1z, z) $$
        \vspace{-0.15cm}
        
        \begin{tikzpicture}[scale=0.55]
            \begin{axis}[width=5.5cm, height=5cm, axis lines=middle, ymin=-2, ymax=5]
                \addplot[green!60!black, thick, domain=-5:5, samples=100] {max(0.1*x,x)};
            \end{axis}
        \end{tikzpicture}
        
        \vspace{0.05cm}
        
        \footnotesize{Prevents "dead neurons".}
    \end{columns}
\end{frame}

\begin{frame}{More Activation Functions}
    \begin{columns}[T]
        \column{0.24\textwidth}
        \centering
        \textbf{Soft-thresholding}
        \vspace{-0.35cm}
        
        $$ \mathcal{S}_\alpha(z) $$
        \vspace{-0.15cm}
        
        \begin{tikzpicture}[scale=0.55]
            \begin{axis}[width=5.5cm, height=5cm, axis lines=middle, xmin=-4, xmax=4, ymin=-3.5, ymax=3.5]
                \addplot[blue, thick, domain=-4:-1] {x+1};
                \addplot[blue, thick, domain=-1:1] {0};
                \addplot[blue, thick, domain=1:4] {x-1};
            \end{axis}
        \end{tikzpicture}
        
        \vspace{0.05cm}
        
        \footnotesize{Shrinks values towards zero, used in sparse recovery.}

        \column{0.24\textwidth}
        \centering
        \textbf{PReLU}
        \vspace{-0.35cm}
        
        $$ \max(\alpha z, z) $$
        \vspace{-0.15cm}
        
        \begin{tikzpicture}[scale=0.55]
            \begin{axis}[width=5.5cm, height=5cm, axis lines=middle, xmin=-4, xmax=4, ymin=-2, ymax=4]
                \addplot[red, thick, domain=-4:0] {0.25*x};
                \addplot[red, thick, domain=0:4] {x};
            \end{axis}
        \end{tikzpicture}
        
        \vspace{0.05cm}
        
        \footnotesize{Similar to Leaky ReLU, but $\alpha$ is a learned parameter.}

        \column{0.24\textwidth}
        \centering
        \textbf{Tanh}
        \vspace{-0.35cm}
        
        $$ \tanh(z) $$
        \vspace{-0.15cm}
        
        \begin{tikzpicture}[scale=0.55]
            \begin{axis}[width=5.5cm, height=5cm, axis lines=middle, xmin=-4, xmax=4, ymin=-1.5, ymax=1.5]
                \addplot[green!60!black, thick, domain=-4:4, samples=100] {tanh(x)};
            \end{axis}
        \end{tikzpicture}
        
        \vspace{0.05cm}
        
        \footnotesize{Zero-centred output, smooth non-linearity.}

        \column{0.24\textwidth}
        \centering
        \textbf{Softmax (2D)}
        \vspace{-0.35cm}
        
        $$ \frac{e^{z_1}}{e^{z_1}+e^{z_2}} $$
        \vspace{-0.15cm}
        
        \begin{tikzpicture}[scale=0.55]
            \begin{axis}[width=5.5cm, height=5cm, view={-30}{30}, colormap/viridis, zmin=-0.1, zmax=1.1, xlabel={$z_1$}, ylabel={$z_2$}, xtick=\empty, ytick=\empty, ztick=\empty]
                \addplot3[surf, domain=-3:3, domain y=-3:3, samples=20, shader=faceted] {exp(x)/(exp(x)+exp(y))};
            \end{axis}
        \end{tikzpicture}
        
        \vspace{0.05cm}
        
        \footnotesize{Maps a vector into a probability distribution.}
    \end{columns}

    \vspace{0.15cm}
    \centering
    And many more...
\end{frame}

\begin{frame}{The Perceptron and the XOR Problem}
    The single layer Perceptron is a \textbf{linear classifier}. It separates data with a hyperplane $w^\top x + b = 0$.
    
    \begin{alertblock}{Limitations \citep{minsky1969}}
        A single perceptron cannot solve the \textbf{XOR problem} because the classes are not linearly separable.
    \end{alertblock}
    
    \centering
    \begin{tikzpicture}[scale=0.8]
        \draw[->] (-0.5,0) -- (2.5,0) node[right] {$x_1$};
        \draw[->] (0,-0.5) -- (0,2.5) node[above] {$x_2$};
        
        % XOR Data points
        \node[draw, circle, fill=red, inner sep=2pt] at (0,0) {}; % 0,0 -> 0
        \node[draw, circle, fill=red, inner sep=2pt] at (2,2) {}; % 1,1 -> 0
        \node[draw, cross out, draw=blue, thick, inner sep=2pt] at (0,2) {}; % 0,1 -> 1
        \node[draw, cross out, draw=blue, thick, inner sep=2pt] at (2,0) {}; % 1,0 -> 1
        
        \node at (1, -1) {No single line can separate red circles from blue crosses.};
    \end{tikzpicture}
\end{frame}

\begin{frame}{Shallow Neural Networks}
    \only<1-2>{To solve XOR, we add a \textbf{hidden layer}.\vspace{-0.1cm}}
    
    \begin{definition}[Shallow Network]
        $$ f_w(x) = \sum_{j = 1}^J c_j \, \sigma(w_j^\top \, x + b_j) $$
    \end{definition}\pause
    
    \vspace{-0.1cm}
        \begin{theoremblock}{Universal Approximation Theorem \protect\citep{cybenko1989,hornik1989,hornik1991}}
        A shallow network with a sufficiently large (but finite) number of neurons $J$ can approximate any continuous function on a compact set to arbitrary precision.
    \end{theoremblock}\pause
    
    \vspace{-0.1cm}
    \only<3->{\textbf{Why Deep?} Shallow networks are inefficient. Representing complex functions might require exponentially many neurons in one layer. Depth allows for \textbf{hierarchical features}.}
\end{frame}

\begin{frame}{Multi-Layer Perceptrons (MLP)}
    Mathematically, a Multi-Layer Perceptron (MLP) or feed-forward neural network \citep{goodfellow2016} is a \textbf{composition} of parameterised functions, i.e.,
    \begin{align*}
        f_w(x) := \varphi_L \circ \varphi_{L-1} \circ \cdots \circ \varphi_1(x)
    \end{align*}
    where each layer is typically affine-linear followed by a non-linearity
    $$ \varphi_l(z) = \sigma(W_l z + b_l) \, . $$
    \begin{center}
    \begin{tikzpicture}[scale=0.7, transform shape]
        % Input Layer
        \foreach \i in {1,...,3}
            \node[circle, draw, fill=green!20] (I\i) at (0,-\i) {$x_\i$};
            
        % Hidden Layer 1
        \foreach \h in {1,...,4}
            \node[circle, draw, fill=blue!20] (H1\h) at (3,-\h+0.5) {};
            
        % Hidden Layer 2
        \foreach \h in {1,...,4}
            \node[circle, draw, fill=blue!20] (H2\h) at (6,-\h+0.5) {};
            
        % Output Layer
        \node[circle, draw, fill=red!20] (O1) at (9,-2) {$y$};
        
        % Connections
        \foreach \i in {1,...,3}
            \foreach \h in {1,...,4}
                \draw[->, opacity=0.3] (I\i) -- (H1\h);
                
        \foreach \h in {1,...,4}
            \foreach \k in {1,...,4}
                \draw[->, opacity=0.3] (H1\h) -- (H2\k);
                
        \foreach \k in {1,...,4}
            \draw[->, opacity=0.3] (H2\k) -- (O1);
            
        \node at (0, -4) {Input};
        \node at (4.5, -4) {Hidden Layers (Feature Extraction)};
        \node at (9, -4) {Output};
    \end{tikzpicture}
    \end{center}
\end{frame}

%================================================================%
\section{Convolutional Neural Networks (CNNs)}
%================================================================%

\begin{frame}{From MLPs to CNNs: Motivation and notation}
    \begin{itemize}
        \item \textbf{Problem:} Applying standard MLPs to high-resolution images results in massive weight matrices.
        \item A $256 \times 256$ RGB image has nearly $200,000$ input values. A single fully connected hidden layer with $1000$ neurons would require $\sim 200,000,000$ weights!
        \item Furthermore, MLPs ignore the inherent two- or three-dimensional spatial structure of images.
    \end{itemize}
    \vspace{0.5em}
    \textbf{Notation:} Let $x \in \mathbb{R}^{H \times W \times C}$ denote an image of height $H$, width $W$, and $C$ channels (for example, $C=3$ for RGB).
\end{frame}

\begin{frame}{From MLPs to CNNs: what changes?}
    \textbf{Connection to MLPs:}
    \begin{itemize}
        \item In an MLP, each layer computes $z = W x + b$ with $W \in \mathbb{R}^{d_{\text{out}} \times d_{\text{in}}}$, treating the input as a vector and ignoring spatial structure.
        \item For images, this means flattening $x$ to a vector, losing locality and spatial relationships.
        \item \textbf{CNNs replace the dense matrix multiplication $W x$ with a convolution operation $x * k$ that preserves spatial structure and uses far fewer parameters.}
    \end{itemize}
    \vspace{0.5em}
    This shift enables the network to exploit local patterns and translation equivariance, which are essential for image data.
\end{frame}

\begin{frame}{Discrete convolutions}
    Instead of every output neuron looking at every input pixel, each neuron only looks at small local patch (e.g., $3 \times 3$) and the same weights (the \textbf{kernel} or \textbf{filter}) are slid across the image.
    
    \textbf{What changes from MLPs?}
    \begin{itemize}
        \item \textbf{Dense layer:} $z = W x + b$ (fully connected, ignores spatial structure)
        \item \textbf{Convolutional layer:} $y = x * k + b$ (local, spatially structured, weight sharing)
        \item The convolution replaces the dense matrix multiplication, drastically reducing parameters and preserving spatial relationships.
    \end{itemize}
    
    \textbf{Mathematical Definition:} for a single-channel 2D image $x \in \mathbb{R}^{H \times W}$ and kernel $k \in \mathbb{R}^{m \times n}$, the discrete convolution is defined as
    \begin{align*}
        (x * k)[i,j] := \sum_{u=0}^{m-1} \sum_{v=0}^{n-1} k[u,v] \; x[i-u, j-v] \, ,
    \end{align*}
    where $[i,j]$ indexes the output location, and $k[u,v]$ are the kernel weights.

\end{frame}


\begin{frame}{Discrete convolutions}
    \textbf{Multi-channel case:} For $C$ input channels and $C'$ output channels, the convolutional layer computes
    \begin{align*}
        y_{c'}[i,j] := \sum_{c=1}^C \sum_{u=0}^{m-1} \sum_{v=0}^{n-1} k_{c',c}[u,v] \; x_c[i-u, j-v] + b_{c'}
    \end{align*}
    where $k_{c',c}$ is the kernel for output channel $c'$ and input channel $c$, and $b_{c'}$ is a bias term.

    \textbf{Benefits:}
    \begin{itemize}
        \item Drastic reduction in parameter count: $m \times n \times C \times C'$ instead of $HWC \times C'$.
        \item \textbf{Translation equivariance:} The same kernel is applied at every spatial location.
        \item Discrete convolution is a linear operator, and a special case of a matrix $W_l$ in MLPs.
    \end{itemize}
\end{frame}

\begin{frame}{Convolution vs. Cross-Correlation in Practice}
    \textbf{Theory vs. Implementation:}
    \begin{itemize}
        \item The strict mathematical definition of convolution uses subtraction ($i-u, j-v$), which implicitly ``flips'' the kernel horizontally and vertically.
        \item Deep learning frameworks (e.g., PyTorch, TensorFlow) actually implement \textbf{discrete cross-correlation}, which uses addition ($i+u, j+v$) and does not flip the kernel.
    \end{itemize}
    
    \vspace{0.5em}
    \textbf{Why omit the flip?}
    \begin{itemize}
        \item Cross-correlation is mathematically equivalent to a convolution with a flipped (transposed) kernel.
        \item Since the kernel weights $k$ are initialized randomly and learned from the data during training, learning a flipped kernel is functionally identical to learning the original kernel. 
        \item Omitting the flip operation simplifies the implementation and the mental model (a sliding window computing dot products) without sacrificing any representational power.
    \end{itemize}
\end{frame}

\begin{frame}{Pooling}
    \textbf{Pooling:} Downsampling operation to reduce spatial dimensions and introduce invariance.
    
    \vspace{0.5em}
    \textbf{What changes from MLPs?}
    \begin{itemize}
        \item In MLPs, dimensionality reduction is achieved by reducing the number of neurons in a layer (via $W$).
        \item In CNNs, \textbf{pooling layers} (e.g., max or average pooling) explicitly reduce spatial resolution while preserving local features and introducing invariances.
        \item Pooling replaces the role of downsampling performed by dense layers, but does so in a spatially aware, local fashion.
    \end{itemize}
    
    \textbf{Example: Max Pooling:}
    \begin{align*}
        y[i,j] = \max_{(u,v) \in \mathcal{P}} x[i+u, j+v]
    \end{align*}
    where $\mathcal{P}$ is the pooling window (e.g., $2 \times 2$).
\end{frame}

\begin{frame}{Pooling, Subsampling, and Receptive Fields}
\textbf{Max Pooling:}
    \begin{align*}
        y[i,j] = \max_{(u,v) \in \mathcal{P}} x[i+u, j+v]
    \end{align*}

    \textbf{Average Pooling:}
    \begin{align*}
        y[i,j] = \frac{1}{|\mathcal{P}|} \sum_{(u,v) \in \mathcal{P}} x[i+u, j+v]
    \end{align*}
    where $\mathcal{P}$ is the pooling window (e.g., $2 \times 2$).
    
    \only<2>{\textbf{Purpose:}
    \begin{itemize}
        \item Dimensionality reduction
        \item Introduce spatial invariance
        \item Extract dominant features
    \end{itemize}}
    
    \only<3->{\textbf{Receptive Field:}
    The region of the input space that affects a particular output unit. Stacking convolutions and pooling increases the receptive field size.}
\end{frame}

%================================================================%
\section{Training Dynamics \& Optimisation}
%================================================================%

\begin{frame}{Empirical Risk Minimisation}
    We train the network by minimising an empirical risk over a dataset $\{(x_i, y_i)\}_{i=1}^s$, i.e.,
    
    $$ \min_{w} \mathcal{L}(w) := \frac{1}{s} \sum_{i=1}^s \ell(f_w(x_i), y_i) \, . $$
    Here $\ell$ denotes a loss function that measures the deviation between the neural network output $f_w(x_i)$ with input $x_i$ and the target $y_i$.\pause

    \vspace{0.5em}
    
    \textbf{Example Loss Functions:}
    \begin{itemize}
        \item Mean Squared Error (MSE): $\frac{1}{2}\| \hat{y} - y \|^2$ (Regression / Inverse Problems).
        \item Cross-Entropy Loss: $-\sum y_k \log(\hat{y}_k)$ (Classification).
    \end{itemize}

    \vspace{0.5em}
    
    \textbf{The Challenge:} The loss landscape $\mathcal{L}(w)$ is highly \textbf{non-convex}.
\end{frame}

\begin{frame}{Convexity}
    \begin{columns}[T]
        \column{0.55\textwidth}
        \begin{block}{Definition}
            A differentiable function $f \colon \mathbb{R}^n \to \mathbb{R}$ is \textbf{convex} if, for all $u, v \in \mathbb{R}^n$ and all $\lambda \in [0,1]$,
            \begin{align*}
                f\!\left( \lambda u + (1-\lambda) v \right) \leq \lambda f(u) + (1-\lambda) f(v) \, .
            \end{align*}
        \end{block}

        \vspace{0.5em}
        Equivalently, $f$ lies \emph{above} every tangent hyperplane, i.e.,
        \begin{align*}
            f(v) \geq f(u) + \left\langle \nabla f(u),\, v - u \right\rangle \, ,
        \end{align*}
        for all $u, v \in \mathbb{R}^n$.

        \column{0.45\textwidth}
        \centering
        \begin{tikzpicture}[scale=0.9]
            % Pre-compute all values
            \pgfmathsetmacro{\ua}{0.7}
            \pgfmathsetmacro{\va}{2.6}
            \pgfmathsetmacro{\fu}{\ua*\ua - \ua + 0.5}
            \pgfmathsetmacro{\fv}{\va*\va - \va + 0.5}
            \pgfmathsetmacro{\slope}{(\fv-\fu)/(\va-\ua)}
            \pgfmathsetmacro{\dfu}{2*\ua - 1}
            \pgfmathsetmacro{\mx}{0.5*\ua + 0.5*\va}
            \pgfmathsetmacro{\mychord}{0.5*\fu + 0.5*\fv}
            \pgfmathsetmacro{\myfunc}{\mx*\mx - \mx + 0.5}
            \pgfmathsetmacro{\tangentend}{\dfu*(3.1-\ua)+\fu}
            \pgfmathsetmacro{\chordlabely}{\fv+0.3}
            \begin{axis}[
                width=6cm, height=5.5cm,
                axis lines=middle,
                xmin=-0.3, xmax=3.3,
                ymin=-0.5, ymax=5.5,
                xtick=\empty, ytick=\empty,
                xlabel={$w$}, ylabel={$f(w)$},
                xlabel style={right},
                ylabel style={above},
                clip=false,
            ]
            % Convex function: f(x) = x^2 - x + 0.5
            \addplot[uclblue, thick, domain=0.1:3.1, samples=80] {x^2 - x + 0.5};
            % Chord between u and v
            \addplot[darkgreen, thick, dashed, domain=0.7:2.6] {\slope*(x-\ua)+\fu};
            % Tangent at u
            \addplot[red, thick, domain=0.1:3.1] {\dfu*(x-\ua)+\fu};
            % Markers at u and v
            \addplot[mark=*, mark size=2pt, uclblue] coordinates {(\ua,\fu)};
            \addplot[mark=*, mark size=2pt, uclblue] coordinates {(\va,\fv)};
            % Midpoint: chord above function
            \addplot[mark=*, mark size=2pt, darkgreen] coordinates {(\mx,\mychord)};
            \addplot[mark=*, mark size=2pt, red!70!black] coordinates {(\mx,\myfunc)};
            \draw[<->, gray, thin] (axis cs:\mx,\myfunc) -- (axis cs:\mx,\mychord)
                node[midway, left, font=\tiny] {gap};
            \node[below, font=\scriptsize] at (axis cs:\ua,\fu) {$u$};
            \node[below, font=\scriptsize] at (axis cs:\va,\fv) {$v$};
            \node[right, font=\scriptsize, red] at (axis cs:3.1,\tangentend) {tangent};
            \node[right, font=\scriptsize, darkgreen] at (axis cs:2.6,\chordlabely) {chord};
            \end{axis}
        \end{tikzpicture}
        \vspace{0.1em}

        {\footnotesize Example: $f(w) = w^2 - w + 0.5$. The function (blue) lies below the chord (green) and above the tangent (red).}
    \end{columns}
\end{frame}

\begin{frame}{$L$-smoothness}
    \begin{columns}[T]
        \column{0.55\textwidth}
        \begin{block}{Definition}
            A differentiable function $f$ is \textbf{$L$-smooth} if there exists $L > 0$ such that, for all $u, v \in \mathbb{R}^n$,
            \begin{align*}
                \left\| \nabla f(u) - \nabla f(v) \right\| \leq L \left\| u - v \right\| .
            \end{align*}
        \end{block}
        Equivalently, $f$ lies \emph{below} the quadratic upper bound at every point,\vspace{-0.25cm}
        \begin{align*}
            f(v) \leq f(u) + \left\langle \nabla f(u),\, v - u \right\rangle + \frac{L}{2}\left\| v - u \right\|^2 \, .
        \end{align*}

        \vspace{-0.25cm}
        The constant $L$ controls how fast the gradient can change: large $L$ means highly curved.

        \vspace{0.25em}
        For example, $f(w) = 1 - \cos\left( w-u \right)$ is globally $1$-smooth because
        \begin{align*}
            f''(w) = \cos\left( w-u \right) \qquad \text{and hence} \qquad \left| f''(w) \right| \leq 1 \, .
        \end{align*}

        \column{0.45\textwidth}
        \centering
        \begin{tikzpicture}[scale=0.9]
            % Pre-compute all values
            \pgfmathsetmacro{\ub}{1.2}
            \pgfmathsetmacro{\Lsmooth}{1.0}
            \pgfmathsetmacro{\fub}{0.0}
            \pgfmathsetmacro{\dfub}{0.0}
            \pgfmathsetmacro{\quadend}{\fub + 0.5*\Lsmooth*(3.0-\ub)*(3.0-\ub) + 0.1}
            \pgfmathsetmacro{\vx}{2.75}
            \pgfmathsetmacro{\dx}{\vx-\ub}
            \pgfmathsetmacro{\fvis}{1 - cos(deg(\dx))}
            \pgfmathsetmacro{\qvis}{\fub + 0.5*\Lsmooth*\dx*\dx}
            \begin{axis}[
                width=6cm, height=5.5cm,
                axis lines=middle,
                xmin=-0.3, xmax=3.3,
                ymin=-0.2, ymax=2.2,
                xtick=\empty, ytick=\empty,
                xlabel={$w$}, ylabel={$f(w)$},
                xlabel style={right},
                ylabel style={above},
                clip=false,
            ]
            % L-smooth function with exact global bound L = 1
            \addplot[uclblue, thick, domain=0.1:3.1, samples=120]
                {1 - cos(deg(x-\ub))};
            % Quadratic upper bound at u
            \addplot[darkgreen, thick, dashed, domain=0.1:3.1, samples=120]
                {\fub + \dfub*(x-\ub) + 0.5*\Lsmooth*(x-\ub)^2};
            % Marker at u
            \addplot[mark=*, mark size=2pt, uclblue] coordinates {(\ub,\fub)};
            \draw[<->, gray, thin] (axis cs:\vx,\fvis) -- (axis cs:\vx,\qvis)
                node[midway, right, font=\tiny] {gap};
            \node[below left, font=\scriptsize] at (axis cs:\ub,\fub) {$u$};
            \node[right, font=\scriptsize, darkgreen] at (axis cs:3.0,\quadend)
                {$f(u){+}\langle\nabla f(u),{\cdot}\rangle{+}\tfrac{L}{2}\|{\cdot}\|^2$};
            \end{axis}
        \end{tikzpicture}

        {\footnotesize Example: $f(w) = 1 - \cos\left( w-u \right)$ with $u = 1.2$. Here the exact global smoothness constant is $L = 1$, and the function (blue) stays below the quadratic envelope (green).}
    \end{columns}
\end{frame}

\begin{frame}{Gradient descent}
    To minimise the empirical risk $\mathcal{L}(w)$, the simplest first-order method is \textbf{gradient descent}.
    
    \textbf{First-order optimality condition:} a minimiser $w^*$ satisfies $\nabla_w \mathcal{L}(w^*) = 0$.\pause

    \vspace{0.5em}
    Starting from an initial guess $w^0$, each iteration performs the update
    \begin{align*}
        w^{k+1} = w^k - \eta \, \nabla_w \mathcal{L}(w^k) \, ,
    \end{align*}
    where $\eta > 0$ is the \textbf{step size} (or learning rate), chosen using smoothness assumptions of $\mathcal{L}$.\pause

    \begin{block}{Convergence}
        If $\mathcal{L}$ is convex and $\nabla_w \mathcal{L}$ is $L$-Lipschitz continuous, then gradient descent converges for any step size $\eta \leq 1/L$, with rate $\mathcal{O}(1/k)$.
    \end{block}

    \textbf{Caveat:} the empirical risk $\mathcal{L}(w)$ in deep learning is usually highly non-convex, so convergence guarantees to global minima do not apply, but in practice, gradient descent often still works remarkably well.
\end{frame}

\begin{frame}{Bregman proximal view of gradient descent}
    Gradient descent has a deeper variational interpretation via the \textbf{Bregman distance} \citep{bregman, kiwiel1997proximal}.

    The Bregman distance generated by a convex function $J$ is
    \begin{align*}
        D_J(w, v) := J(w) - J(v) - \left\langle \nabla J(v),\, w - v \right\rangle \geq 0 \, .
    \end{align*}
    
    \vspace{-0.15cm}
    \pause

    \begin{columns}[T]
        \column{0.47\textwidth}
        Consider the convex function $J(w) = w \log\left( w \right) - w$ on $\left( 0, \infty \right)$, for which
        \begin{align*}
            \nabla J(w) = \log\left( w \right) .
        \end{align*}
        Hence the Bregman distance is the vertical gap
        \begin{align*}
            D_J(w, v) = w \log\left( \frac{w}{v} \right) - w + v \, .
        \end{align*}

        \column{0.51\textwidth}
        \centering
        \begin{tikzpicture}[scale=0.9]
            % Example parameters
            \pgfmathsetmacro{\vb}{1.3}
            \pgfmathsetmacro{\wb}{2.45}
            \pgfmathsetmacro{\Jv}{\vb*ln(\vb) - \vb}
            \pgfmathsetmacro{\dJv}{ln(\vb)}
            \pgfmathsetmacro{\Jw}{\wb*ln(\wb) - \wb}
            \pgfmathsetmacro{\Tw}{\Jv + \dJv*(\wb-\vb)}
            \pgfmathsetmacro{\tanlabely}{\Jv + \dJv*(3.2-\vb)}
            \begin{axis}[
                width=6.2cm, height=5.5cm,
                axis lines=middle,
                xmin=0.2, xmax=3.4,
                ymin=-1.2, ymax=1.0,
                xtick=\empty, ytick=\empty,
                xlabel={$w$}, ylabel={$J(w)$},
                xlabel style={right},
                ylabel style={above},
                clip=false,
            ]
            \addplot[uclblue, thick, domain=0.25:3.2, samples=140] {x*ln(x) - x};
            \addplot[red, thick, domain=0.25:3.2, samples=120] {\Jv + \dJv*(x-\vb)};
            \addplot[mark=*, mark size=2pt, uclblue] coordinates {(\vb,\Jv)};
            \addplot[mark=*, mark size=2pt, uclblue] coordinates {(\wb,\Jw)};
            \addplot[mark=*, mark size=2pt, red] coordinates {(\wb,\Tw)};
            \draw[<->, gray, thin] (axis cs:\wb,\Tw) -- (axis cs:\wb,\Jw)
                node[midway, right, font=\tiny] {$D_J(w,v)$};
            \node[below, font=\scriptsize] at (axis cs:\vb,\Jv) {$v$};
            \node[below, font=\scriptsize] at (axis cs:\wb,\Tw) {$w$};
            \node[right, font=\scriptsize, red] at (axis cs:3.2,\tanlabely) {tangent at $v$};
            \end{axis}
        \end{tikzpicture}

        % {\footnotesize Example: $J(w) = w\log\left( w \right) - w$. The Bregman distance is the gap between $J$ and the tangent line at $v$.}
    \end{columns}
\end{frame}

\begin{frame}{Bregman proximal view of gradient descent}
    Gradient descent has a deeper variational interpretation via the \textbf{Bregman distance}.

    The Bregman distance generated by a convex function $J$ is
    \begin{align*}
        D_J(w, v) := J(w) - J(v) - \left\langle \nabla J(v),\, w - v \right\rangle \geq 0 \, .
    \end{align*}
    \vspace{0.5em}
    Consider the Bregman proximal iteration
    \begin{align*}
        w^{k+1} = \argmin_{w} \left\{ \mathcal{L}(w) + D_J(w, w^k) \right\}
    \end{align*}
    Different choices of $J$ produce different local geometries. Gradient descent arises from one particular choice.
\end{frame}

\begin{frame}{Bregman proximal view of gradient descent (cont.)}
    For gradient descent, we choose
    \begin{align*}
        J(w) = \frac{1}{2\eta}\|w\|^2 - \mathcal{L}(w) ,
    \end{align*}
    so that $\nabla J(w) = \frac{1}{\eta}w - \nabla_w \mathcal{L}(w)$. Expanding the Bregman distance then yields
    \begin{align*}
        \mathcal{L}(w) + D_J(w, w^k)
        &= \mathcal{L}(w^k) + \left\langle \nabla_w \mathcal{L}(w^k),\, w - w^k \right\rangle + \frac{1}{2\eta}\left\| w - w^k \right\|^2 \, .
    \end{align*}
    This is precisely the \textbf{local quadratic model} of $\mathcal{L}$ at $w^k$.

    \vspace{0.5em}
    Taking the first-order optimality condition of this model yields
    \begin{align*}
        0 = \frac{1}{\eta}\left( w^{k+1} - w^k \right) + \nabla_w \mathcal{L}(w^k)
        \quad \Longleftrightarrow \quad
        w^{k+1} = w^k - \eta \, \nabla_w \mathcal{L}(w^k) \, ,
    \end{align*}
    recovering the gradient descent update.

    \vspace{0.5em}
    \begin{theoremblock}{Key insight}
        Gradient descent is a Bregman proximal step. The choice of $J$ determines the geometry of the descent, and other choices of $J$ lead to natural generalisations (mirror descent, natural gradient).
    \end{theoremblock}
\end{frame}

\begin{frame}{Stochastic gradient method (SGM)}
    \textbf{Problem:} computing $\nabla_w \mathcal{L}(w^k) = \frac{1}{s}\sum_{i=1}^s \nabla_w \ell(f_{w^k}(x_i), y_i)$ requires a pass over the entire dataset at each iteration; often infeasible for large $s$.\pause

    \vspace{0.5em}
    \textbf{Key observation:} each per-sample gradient is an \emph{unbiased estimator} of the full gradient,
    \begin{align*}
        \mathbb{E}_i\!\left[ \nabla_w \ell\!\left( f_w(x_i), y_i \right) \right]
        = \frac{1}{s}\sum_{j=1}^s \nabla_w \ell\!\left( f_w(x_j), y_j \right)
        = \nabla_w \mathcal{L}(w) \, .
    \end{align*}
    
    \pause

    \textbf{SGM update:} at each iteration, draw an index $i \sim \mathcal{U}\{1,\ldots,s\}$ uniformly at random and set
    \begin{align*}
        w^{k+1} = w^k - \eta_k \, \nabla_w \ell\!\left( f_{w^k}(x_i), y_i \right) .
    \end{align*}
    
    \pause

    \begin{alertblock}{SGM is not a descent method}
        There is no guarantee that $\mathcal{L}(w^{k+1}) \leq \mathcal{L}(w^k)$ at every step. The injected stochastic noise can, however, help escape sharp local minima.
    \end{alertblock}
\end{frame}

\begin{frame}{Mini-batching and step size schedules}
    \textbf{Mini-batch SGM:} instead of a single sample, compute the gradient over a random subset $\mathcal{B}_k \subset \{1,\ldots,s\}$ of batch size $B = |\mathcal{B}_k| \ll s$,
    \begin{align*}
        w^{k+1} = w^k - \eta_k \, \frac{1}{B} \sum_{i \in \mathcal{B}_k} \nabla_w \ell\!\left( f_{w^k}(x_i), y_i \right) .
    \end{align*}
    This reduces gradient variance whilst remaining cheaper per iteration than gradient descent.\pause

    \textbf{Step size schedules:} convergence of SGM (in the convex setting) requires a diminishing sequence $\{\eta_k\}$ satisfying the Robbins--Monro conditions, i.e.,
    \begin{align*}
        \lim_{k \to \infty} \eta_k = 0 \qquad \text{and} \qquad \sum_{k=0}^\infty \eta_k = \infty \, .
    \end{align*}
    \only<1-2>{Typical choices are $\eta_k = \eta/(k+1)$ or $\eta_k = \eta/\sqrt{k+1}$. Under these conditions, SGM converges at rate $\mathcal{O}(1/\sqrt{k})$ — slower than the $\mathcal{O}(1/k)$ rate of gradient descent.}

    \only<3->{In practice, fixed or piece-wise constant step sizes with warm restarts are used with early stopping, often outperforming theoretically motivated schedules.}
\end{frame}

\begin{frame}{From momentum to Adam}
    Let $g^k := \nabla_w \ell\!\left( f_{w^k}(x_{i_k}), y_{i_k} \right)$ denote a stochastic gradient.

    \vspace{0.5em}
    \textbf{Momentum (first moment smoothing).}
    \begin{align*}
        m^k &= \beta_1 m^{k-1} + (1-\beta_1) g^k \, , \\
        w^{k+1} &= w^k - \eta \, m^k \, .
    \end{align*}
    Here $m^k$ is an exponential moving average of the past gradients, i.e.,
    \begin{align*}
        m^k = (1-\beta_1)\sum_{t=1}^{k} \beta_1^{k-t} g^t \, .
    \end{align*}

    \vspace{0.25em}
    \textbf{Root Mean Square (RMS) scaling (second moment smoothing).}
    \begin{align*}
        v^k &= \beta_2 v^{k-1} + (1-\beta_2)\left( g^k \odot g^k \right) \, .
    \end{align*}
    This estimates gradient magnitudes and enables step-size normalisation \textbf{per-coordinate}.
\end{frame}

\begin{frame}{Adam: update rule and bias correction}
    Adam combines both moving averages with diagonal preconditioning. Starting with $m^0 = 0$ and $v^0 = 0$, define
    \begin{align*}
        m^k &= \beta_1 m^{k-1} + (1-\beta_1) g^k \, , \\
        v^k &= \beta_2 v^{k-1} + (1-\beta_2)\left( g^k \odot g^k \right) \, .
    \end{align*}
    Since both are initialised at zero, they are biased towards zero for small $k$. Adam therefore uses
    \begin{align*}
        \hat m^k &= \frac{m^k}{1-\beta_1^k} \, ,
        &\hat v^k &= \frac{v^k}{1-\beta_2^k} \, .
    \end{align*}
    The parameter update is\vspace{-0.20cm}
    \begin{align*}
        w^{k+1} = w^k - \eta \, \frac{\hat m^k}{\sqrt{\hat v^k} + \varepsilon} \, ,
    \end{align*}

    \vspace{-0.13cm}
    where the division and square root are applied component-wise.
\end{frame}

\begin{frame}{Adam: interpretation and practical caveats}
    Adam can be written as a diagonally preconditioned momentum step, i.e.,
    \begin{align*}
        w^{k+1} = w^k - \eta \, P_k \, \hat m^k \, ,
        \qquad
        P_k := \operatorname{diag}\!\left( \left( \sqrt{\hat v^k} + \varepsilon \right)^{-1} \right) .
    \end{align*}
    Coordinates with persistently large gradients receive smaller effective steps, while flat coordinates receive larger steps.

    \vspace{0.5em}
    Typical defaults are $\beta_1 = 0.9$, $\beta_2 = 0.999$, and $\varepsilon = 10^{-8}$.

    \vspace{0.5em}
    \begin{alertblock}{Theory vs practice}
        Adam is highly effective in deep learning practice, but classical convergence guarantees can fail without additional assumptions or modifications. A common robust variant is AdamW, which adds weight decay (Tikhonov regularisation) to the objective.
    \end{alertblock}
\end{frame}

\begin{frame}{Practical training of deep networks}
    Successful training of modern deep networks requires more than an optimiser. We must address two fundamental challenges:
    \begin{enumerate}
        \item \textbf{Overfitting and generalisation:} The model memorises training data instead of learning general patterns. Deep networks have millions of parameters, but often only thousands or tens of thousands of training examples.
        \item \textbf{Internal covariate shift:} As weights update, the distribution of each layer's input changes, slowing convergence and requiring careful tuning of step sizes.
    \end{enumerate}
    \vspace{0.5em}
    Two practical techniques address these concerns:
    \begin{itemize}
        \item \textbf{Dropout:} Stochastic regularisation applied during training to prevent neuron co-adaptation and overfitting.
        \item \textbf{Batch normalisation:} Normalises layer inputs to stabilise training dynamics and reduce internal covariate shift.
    \end{itemize}
\end{frame}

\begin{frame}{Dropout in practice}
    \textbf{Placement:} After a layer's activation, typically at each hidden layer. Each neuron is independently dropped with probability $(1-q)$.

    \vspace{0.25em}
    During training, randomly zero out a fraction $(1-q)$ of activations using a Bernoulli mask, i.e.,
    \begin{align*}
        m_j \sim \mathrm{Bernoulli}(q), \qquad j = 1, \ldots, d,
    \end{align*}
    where $q \in (0,1]$ is the \textbf{keep probability} (typical: $q = 0.5$).

    \vspace{0.25em}
    \textbf{Inverted dropout} (modern standard) rescales activations, i.e.,
    \begin{align*}
        \tilde h = \frac{1}{q}\left( m \odot h \right) .
    \end{align*}
    This preserves expected activation: $\mathbb{E}[\tilde h_j] = h_j$, so no test-time adjustment is needed.

    \vspace{0.25em}
    \textbf{Inference:} Dropout is disabled; deterministically use $\tilde h = h$.
\end{frame}

\begin{frame}{Why dropout works}
    \textbf{Ensemble perspective:} Dropout creates an ensemble of $2^d$ sub-networks (one per binary mask). Inference approximates their average prediction, reducing variance.

    \vspace{0.5em}
    \textbf{Regularisation perspective:} The stochastic perturbations encourage the network to learn redundant, robust features and prevent co-adaptation — where neurons become specialised to specific partners in training data.

    \vspace{0.5em}
    \textbf{Effect on optimisation:} By reducing overfitting, dropout improves generalisation to unseen data. The noise also provides implicit regularisation during SGD, helping escape sharp local minima.
\end{frame}

\begin{frame}{Batch normalisation in practice}
    \only<1>{\textbf{Placement:} Applied to pre-activations $z$ (before the nonlinearity) within a hidden layer. Typical architecture: Dense $\to$ Batch Norm $\to$ Nonlinearity $\to$ Dropout.}

    \vspace{0.25em}
    For each feature/channel, compute mini-batch statistics of the form
    \begin{align*}
        \mu_{\mathcal{B}} &= \frac{1}{|\mathcal{B}|} \sum_{i \in \mathcal{B}} z_i, \\[0.1em]
        \sigma_{\mathcal{B}}^2 &= \frac{1}{|\mathcal{B}|} \sum_{i \in \mathcal{B}} \left( z_i - \mu_{\mathcal{B}} \right)^2 .
    \end{align*}\pause
    Then normalise and apply learnable affine parameters, i.e.,
    \begin{align*}
        \hat z_i &= \frac{z_i - \mu_{\mathcal{B}}}{\sqrt{\sigma_{\mathcal{B}}^2 + \varepsilon}}, \\[0.1em]
        y_i &= \gamma \, \hat z_i + \beta .
    \end{align*}
    Here $\gamma$ (scale) and $\beta$ (shift) are trained via backpropagation.
\end{frame}

\begin{frame}{Batch normalisation: training vs inference}
    \textbf{Training:} Use mini-batch statistics $\mu_{\mathcal{B}}$, $\sigma_{\mathcal{B}}^2$.

    \vspace{0.5em}
    \textbf{Inference:} Mini-batches are unavailable (single sample or fixed batch). Instead, use \emph{running averages} accumulated during training via exponential moving average (EMA), i.e.,
    \begin{align*}
        \mu_{\text{running}} &\gets (1-\alpha) \mu_{\text{running}} + \alpha \, \mu_{\mathcal{B}}, \\[0.1em]
        \sigma^2_{\text{running}} &\gets (1-\alpha) \sigma^2_{\text{running}} + \alpha \, \sigma_{\mathcal{B}}^2,
    \end{align*}
    where $\alpha$ is a momentum coefficient (typical: $\alpha = 0.1$).
\end{frame}

\begin{frame}{Why batch normalisation helps}
    \textbf{Internal Covariate Shift Problem:} As weights $w$ update, the distribution of inputs to each layer shifts. Downstream layers must continuously re-adapt, slowing convergence.

    \vspace{0.5em}
    \textbf{Batch Normalisation Solution:} Enforce zero mean and unit variance on layer inputs. Benefits:
    \begin{itemize}
        \item \textbf{Faster convergence:} Reduced covariate shift enables higher learning rates and more stable training dynamics.
        \item \textbf{Mitigates vanishing gradients:} Normalisation keeps pre-activations in the active range of nonlinearities (e.g., $z \in [-2,2]$ for tanh), preventing gradient attenuation through deep layers.
        \item \textbf{Implicit regularisation:} Mini-batch noise acts as a regulariser, reducing overfitting similar to dropout.
        \item \textbf{Learnable scale/shift:} Parameters $\gamma, \beta$ allow the network to learn the optimal input distribution, preserving expressiveness.
    \end{itemize}
\end{frame}

%================================================================%
\section{The Chain Rule of Calculus \& Backpropagation}
%================================================================%

\begin{frame}{How do we compute gradients?}
    In the optimisation section, we introduced updates such as
    \begin{align*}
        w^{k+1} &= w^k - \eta\,\nabla_w\mathcal{L}(w^k), \\
        w^{k+1} &= w^k - \eta_k\,\nabla_w\ell\!\left(f_{w^k}(x_{i_k}),y_{i_k}\right) .
    \end{align*}
    These formulas are only useful if we can compute the gradients efficiently for deep networks.

    \vspace{0.5em}
    \begin{block}{Core question}
        Given $f_w(x)=\varphi_L\circ\cdots\circ\varphi_1(x)$, how do we compute
        $\nabla_w\mathcal{L}(w)$ and stochastic gradients without differentiating each path separately?
    \end{block}

    Backpropagation \citep{rumelhart1986learning} is the answer: a dynamic-programming application of the chain rule that reuses intermediate quantities and computes all parameter gradients layer-by-layer.
\end{frame}

\begin{frame}{Backpropagation as constrained optimisation}
    For the MLP introduced earlier, i.e.,\vspace{-0.15cm}
    \begin{align*}
        f_w(x) = \varphi_L \circ \cdots \circ \varphi_1(x),
        \qquad
        \varphi_l(z)=\sigma_l(W_l z+b_l) \, ,
    \end{align*}
    we collect all trainable parameters as\vspace{-0.15cm}
    \begin{align*}
        w = \{(W_l,b_l)\}_{l=1}^L \, .
    \end{align*}
    For each sample $(x^i,y^i)$, we define layer variables
    \begin{align*}
        z_l^i &= W_l x_{l-1}^i + b_l, \\
        x_l^i &= \sigma_l(z_l^i),
        \qquad l=1,\ldots,L,
        \qquad x_0^i = x^i \, .
    \end{align*}
    Training minimises the empirical risk\vspace{-0.15cm}
    \begin{align*}
        \min_{w} \; \frac{1}{s}\sum_{i=1}^s \ell\!\left(y^i, x_L^i\right),
    \end{align*}
    where the dependence on $w$ is nested through all layers.
\end{frame}

\begin{frame}{Equivalent constrained formulation}
    Let's introduce activation and pre-activation variables as optimisation variables, i.e.,
    \begin{align*}
        \min_{w,\{x_l^i,z_l^i\}} \; &\frac{1}{s}\sum_{i=1}^s \ell\!\left(y^i, x_L^i\right)
    \end{align*}
    subject to
    \begin{align*}
        z_l^i - W_l x_{l-1}^i - b_l &= 0, \\
        x_l^i - \sigma_l(z_l^i) &= 0,
        \qquad l=1,\ldots,L,
        \quad i=1,\ldots,s \, .
    \end{align*}
    This formulation is equivalent to the nested network definition and makes layerwise dependencies explicit.
\end{frame}

\begin{frame}{Lagrangian and optimality conditions}
    We can then introduce so-called Lagrange multipliers $p_l^i$ and $q_l^i$ for the two constraints, and define the Lagrangian
    \begin{align*}
        \mathcal{L} = \frac{1}{s}\sum_{i=1}^s \Bigg[\ell\!\left(y^i,x_L^i\right)
        + \sum_{l=1}^L \left\langle p_l^i,\, z_l^i - W_l x_{l-1}^i - b_l \right\rangle
        + \sum_{l=1}^L \left\langle q_l^i,\, x_l^i - \sigma_l(z_l^i) \right\rangle\Bigg] \, .
    \end{align*}

    \vspace{-0.1cm}
    Stationarity with respect to $x_l^i$ gives the adjoint recursion
    \begin{align*}
        q_L^i &= -\nabla_{x_L}\ell\!\left(y^i,x_L^i\right), \\
        q_l^i &= W_{l+1}^{\top}p_{l+1}^i,
        \qquad l=L-1,\ldots,1 \, .
    \end{align*}
    Stationarity with respect to $z_l^i$ gives\vspace{-0.2cm}
    \begin{align*}
        p_l^i = \left(J\sigma_l(z_l^i)\right)^{\top} q_l^i \, .
    \end{align*}
    Here $J\sigma_l(z)$ denotes the Jacobian matrix of $\sigma_l$ at $z$. For elementwise activations, $J\sigma_l(z)=\operatorname{diag}(\sigma_l'(z))$.
\end{frame}

\begin{frame}{Recovery of classical backpropagation recursions}
    Stationarity with respect to multipliers recovers the forward pass constraints,
    \begin{align*}
        \nabla_{p_l^i}\mathcal{L}=0 &\Longleftrightarrow z_l^i = W_l x_{l-1}^i + b_l, \\
        \nabla_{q_l^i}\mathcal{L}=0 &\Longleftrightarrow x_l^i = \sigma_l(z_l^i) \, .
    \end{align*}

    \vspace{-0.1cm}
    Let's define pre-activation errors $\delta_l^i := -p_l^i$. Combining the previous relations with these yields
    \begin{align*}
        \delta_L^i &= \nabla_{x_L}\ell\!\left(y^i,x_L^i\right) \odot \sigma_L'\!\left(z_L^i\right), \\
        \delta_l^i &= \left(W_{l+1}^{\top}\delta_{l+1}^i\right) \odot \sigma_l'\!\left(z_l^i\right),
        \qquad l=L-1,\ldots,1 \, .
    \end{align*}
    This is the standard backward recursion from the chain rule, obtained here via first-order optimality of the constrained problem.
\end{frame}

\begin{frame}{Gradients with respect to parameters}
    Stationarity of $\mathcal{L}$ with respect to $W_l$ and $b_l$ then gives
    \begin{align*}
        \nabla_{W_l}\mathcal{L}(w)
        &= \frac{1}{s}\sum_{i=1}^s \delta_l^i\left(x_{l-1}^i\right)^{\top}, \\
        \nabla_{b_l}\mathcal{L}(w)
        &= \frac{1}{s}\sum_{i=1}^s \delta_l^i \, .
    \end{align*}
    A backpropagation step computes all gradients via
    \begin{enumerate}
        \item forward pass for $(x_l^i,z_l^i)$,
        \item backward recursion for $\delta_l^i$,
        \item accumulation of $\nabla_{W_l}\mathcal{L}$ and $\nabla_{b_l}\mathcal{L}$.
    \end{enumerate}
\end{frame}

\begin{frame}{Generalising backpropagation: an operator view}
    Classical backpropagation assumes explicit layer maps
    \begin{align*}
        x_l = \sigma_l(W_l x_{l-1}+b_l).
    \end{align*}
    In modern inverse-problem architectures, a ``layer'' is often defined \emph{implicitly} as the solution of an operator equation.

    \begin{block}{Unified template}
        Define a state $z^\star$ by solving\vspace{-0.15cm}
        \begin{align*}
            \mathcal{G}(z^\star,w,x) = 0,
        \end{align*}
        then optimise a loss $\ell(z^\star,y)$ with respect to parameters $w$.
    \end{block}

    Backpropagation generalises to \textbf{implicit differentiation} and \textbf{adjoint equations} for this constraint.
\end{frame}

\begin{frame}{General constrained Lagrangian form}
    Hence, a unified training model is the equality-constrained problem
    \begin{align*}
        \min_{z,w}\; &\ell(z,y) \\
        \text{s.t.}\; &\mathcal{G}(z,w,x)=0.
    \end{align*}
    Here $z$ is the implicit layer state (or trajectory), $w$ are trainable parameters, and $\mathcal{G}$ defines the operator constraint.

    \vspace{0.5em}
    Let's introduce a Lagrange multiplier $\lambda$ and define
    \begin{align*}
        \mathcal{L}(z,w,\lambda)
        :=\ell(z,y)+\langle \lambda,\mathcal{G}(z,w,x)\rangle.
    \end{align*}
    This is the common starting point for implicit, equilibrium, and continuous-depth backpropagation rules.
\end{frame}

\begin{frame}{General stationarity conditions}
    For equality constraints, the first-order stationarity conditions are
    \begin{align*}
        \nabla_{\lambda}\mathcal{L}=0
        &\Longleftrightarrow \mathcal{G}(z,w,x)=0
        \quad \text{(primal feasibility)}, \\
        \nabla_{z}\mathcal{L}=0
        &\Longleftrightarrow \nabla_z\ell(z,y)+\left(\partial_z\mathcal{G}(z,w,x)\right)^\top\lambda=0, \\
        \nabla_{w}\mathcal{L}=0
        &\Longleftrightarrow \left(\partial_w\mathcal{G}(z,w,x)\right)^\top\lambda=0
        \quad \text{(at optimum)}.
    \end{align*}
    In practice, the second line is the \textbf{adjoint equation}; once $\lambda$ is solved, gradients with respect to $w$ follow directly.
\end{frame}



\begin{frame}{Case I: Bilevel or optimisation-defined layers}
    A layer may be the solution map of an inner optimisation problem \citep{holler2018bilevel, crockett2022bilevel}:
    \begin{align*}
        z^\star(w,x) \in \argmin_{z} E(z,w;x),
        \qquad
        \mathcal{L}(w)=\ell(z^\star(w,x),y).
    \end{align*}

    \textbf{Constraint identification:} The inner optimality condition is exactly $\mathcal{G}=0$:
    \begin{align*}
        \mathcal{G}(z,w,x) := \nabla_z E(z,w;x), \qquad \partial_z\mathcal{G} = \nabla_{zz}^2 E, \quad \partial_w\mathcal{G} = \nabla_{wz}^2 E.
    \end{align*}

    The adjoint equation $\nabla_z\ell + (\partial_z\mathcal{G})^\top\lambda=0$ gives $\lambda = -(\nabla_{zz}^2E)^{-\top}\nabla_z\ell$, and the gradient rule $\nabla_w\mathcal{L}=(\partial_w\mathcal{G})^\top\lambda$ yields
    \begin{align*}
        \frac{\partial z^\star}{\partial w} = -(\nabla_{zz}^2E)^{-1}\nabla_{wz}^2E,
    \end{align*}
    giving hypergradients for bilevel learning without unrolling all inner iterations.
\end{frame}

\begin{frame}{Case II: Deep equilibrium and fixed-point layers}
    In DEQ-style models \citep{bai2019deep, gilton2021deep}, the hidden state is defined by a fixed point:
    \begin{align*}
        z^\star = f_w(z^\star,x).
    \end{align*}

    \textbf{Constraint identification:}
    \begin{align*}
        \mathcal{G}(z,w,x) := f_w(z,x)-z, \qquad \partial_z\mathcal{G} = \partial_z f_w - I, \quad \partial_w\mathcal{G} = \partial_w f_w.
    \end{align*}
    The adjoint equation $\nabla_z\ell + (\partial_z\mathcal{G})^\top\lambda=0$ becomes the linear system
    \begin{align*}
        \left(I-\partial_z f_w(z^\star,x)^\top\right)v = \nabla_z\ell(z^\star,y),
    \end{align*}
    and the gradient rule $(\partial_w\mathcal{G})^\top\lambda=0$ gives
    \begin{align*}
        \nabla_w\mathcal{L}(w) = \left(\partial_w f_w(z^\star,x)\right)^\top v.
    \end{align*}
    This gives constant-memory training with respect to effective depth.
\end{frame}

\begin{frame}{Case III: Continuous networks and adjoint equations}
    For continuous networks, the hidden state $z(t)$ evolves along an ODE:
    \begin{align*}
        \dot z(t)=\mathcal{F}(z(t),w(t)),
        \qquad
        z(0)=x \quad \text{(input)}.
    \end{align*}

    \textbf{Constraint identification:} $\mathcal{G}$ is the ODE residual (a functional operator on $[0,T]$):
    \begin{align*}
        \mathcal{G}[z,w](t) := \dot z(t) - \mathcal{F}(z(t),w(t)) = 0.
    \end{align*}
    The Lagrange multiplier $\lambda$ becomes a continuous adjoint state $p(t)$. The stationarity condition $\nabla_z\mathcal{L}=0$ yields the adjoint ODE
    \begin{align*}
        -\dot p(t)
        =\left(\partial_z\mathcal{F}(z(t),w(t))\right)^*p(t)
        +\nabla_z\mathcal{R}(z(t),w(t)),
        \quad
        p(T)=\nabla_z\mathcal{M}(z(T)),
    \end{align*}
    and $\nabla_w\mathcal{L}=0$ gives the parameter gradient. Backpropagation becomes an adjoint boundary-value computation in continuous time.
\end{frame}

\begin{frame}{Why this matters for the rest of the course}
    We now have one conceptual pipeline:
    \begin{enumerate}
        \item Formulate the layer as a constraint $\mathcal{G}(z,w,x)=0$.
        \item Derive the corresponding adjoint/implicit differentiation rule.
        \item Compute parameter gradients from that rule.
    \end{enumerate}

    \vspace{0.5em}
    This unifies
    \begin{itemize}
        \item classical backpropagation for explicit MLP/CNN layers,
        \item bilevel and optimisation-defined layers,
        \item DEQ fixed-point models,
        \item continuous-depth and optimal-control formulations.
    \end{itemize}
\end{frame}



%================================================================%
\section{Advanced Architectures: ResNets and UNets}
%================================================================%

\begin{frame}{Vanishing Gradients \& Stability}
    \begin{itemize}
        \item \textbf{Problem:} In very deep networks, gradients involve products of many weight matrices and activation derivatives.
        $$ \nabla_{x_0} \mathcal{L} \propto \prod_{l=1}^L W_l \cdot \sigma_l'(z_l) $$
            \item \textbf{Result:} If singular values of $W_l$ are $< 1$, gradients decay exponentially to $0$. If $> 1$, they explode. Training becomes extremely difficult for standard architectures $> 20$ layers \citep{bengio1994learning, glorot2010understanding}.
    \end{itemize}
\end{frame}

\begin{frame}{Residual Networks (ResNets)}
    \textbf{Solution:} Introduce a "Skip Connection" \citep{he2016deep}. Instead of learning $x_{l+1} = \mathcal{F}_w(x_l)$, we learn the \textbf{residual}:
    $$ x_{l+1} = x_l + \mathcal{F}_w(x_l) $$
    Here $\mathcal{F}_w$ is a parameterised block (e.g.\ two conv layers with the same notation as Case III).
    
    \begin{center}
    \begin{tikzpicture}[scale=0.8, transform shape, node distance=2cm]
        \node (in) {$x_l$};
        \node[draw, rectangle, minimum width=2cm, minimum height=1cm, right=of in, align=center] (block) {Layers\\$\mathcal{F}_w$};
        \node[circle, draw, right=of block] (sum) {$+$};
        \node[right=of sum] (out) {$x_{l+1}$};
        
        \draw[->] (in) -- (block);
        \draw[->] (block) -- (sum);
        \draw[->] (sum) -- (out);
        
        % Skip connection
        \draw[->] (in) to[out=90,in=90] node[midway, above] {Identity Skip Connection} (sum);
    \end{tikzpicture}
    \end{center}
    
    \textbf{Benefit:} Creates a "gradient super-highway", allowing gradients to flow unimpeded through the identity path.
\end{frame}

\begin{frame}{ResNets as Continuous Networks (Dynamical Systems View)}
    The skip-connection update $x_{l+1} = x_l + h\,\mathcal{F}_w(x_l)$ is Euler's method applied to
    $$ \dot z(t) = \mathcal{F}(z(t), w), \qquad z(0) = x \quad\text{(input)}, $$
    which is exactly the ODE of Case III with step size $h$ and layer index $l \leftrightarrow t$.

    \begin{theoremblock}{Dynamical Systems View \protect\citep{weinan2017proposal, haber2017stable, benning2019deep}}
        A ResNet is a discretised continuous network solving an initial value problem. Learning is finding the vector field $\mathcal{F}$ (parametrised by $w$) that transports input data $z(0)=x$ to the correct output at terminal time $z(T)$.
    \end{theoremblock}

    In the continuous limit, gradients are computed via the adjoint ODE of Case III rather than discrete backpropagation through all layers.
\end{frame}

\begin{frame}{Encoder-Decoder Architectures and UNets}
    Let $S$ be the number of scales. A UNet \citep{ronneberger2015u} defines $f_w : \mathcal{X} \to \mathcal{X}$ via three stages.

    \textbf{Encoder} (contracting path): starting from $x_0 = x$, for $s = 0, \ldots, S-1$,
    \vspace{-0.2cm}
    \begin{align*}
        e_s &= \mathcal{E}_s(x_s), \qquad x_{s+1} = \mathcal{P}_s(e_s),
    \end{align*}
    where $\mathcal{E}_s$ is a convolutional block and $\mathcal{P}_s$ applies $2\times$ spatial downsampling (e.g., max pooling).

    \textbf{Decoder} (expanding path): starting from the bottleneck $d_S = \mathcal{E}_S(x_S)$, for $s = S-1, \ldots, 0$,
    \vspace{-0.2cm}
    \begin{align*}
        d_s = \mathcal{D}_s\!\left(\operatorname{cat}(\mathcal{U}_s(d_{s+1}),\, e_s)\right),
    \end{align*}
    where $\mathcal{U}_s$ applies $2\times$ upsampling, $\mathcal{D}_s$ is a convolutional block, and $\operatorname{cat}$ denotes channel-wise concatenation.

    \textbf{Output}: $f_w(x) = \mathcal{C}(d_0)$ via a final $1\!\times\!1$ convolution $\mathcal{C}$.

    The skip connections $e_s \to d_s$ bypass the bottleneck, preserving high-frequency spatial detail that is otherwise lost under downsampling.
\end{frame}

\begin{frame}{UNet Architecture Diagram}
    \begin{center}
    \resizebox{0.8\textwidth}{!}{
    \begin{tikzpicture}[
        box/.style={draw, thick, fill=uclblue!20, minimum width=1.2cm, align=center},
        conv/.style={->, thick, solid},
        pool/.style={->, thick, red},
        up/.style={->, thick, uclblue},
        skip/.style={->, thick, dashed, darkgreen}
    ]
        % Encoder
        \node[box, minimum height=4cm] (e0) at (0, 0) {$e_0$};
        \node[box, minimum height=3cm] (e1) at (2, -2.5) {$e_1$};
        \node[box, minimum height=2cm] (e2) at (4, -5) {$e_2$};
        
        % Bottleneck
        \node[box, minimum height=1.5cm, fill=red!20] (b) at (6, -7) {$e_3$};
        
        % Decoder
        \node[box, minimum height=2cm] (d2) at (8, -5) {$d_2$};
        \node[box, minimum height=3cm] (d1) at (10, -2.5) {$d_1$};
        \node[box, minimum height=4cm] (d0) at (12, 0) {$d_0$};
        
        % Input / Output
        \node[left=1.5cm of e0] (in) {$x$};
        \node[right=1.5cm of d0] (out) {$f_w(x)$};
        
        \draw[conv] (in) -- node[above] {in} (e0);
        \draw[conv] (d0) -- node[above] {out} (out);
        
        % Pooling (Contracting)
        \draw[pool] (e0.south) |- node[pos=0.75, above, scale=0.8, text=red] {Pool} (e1.west);
        \draw[pool] (e1.south) |- node[pos=0.75, above, scale=0.8, text=red] {Pool} (e2.west);
        \draw[pool] (e2.south) |- node[pos=0.75, above, scale=0.8, text=red] {Pool} (b.west);
        
        % Upsampling (Expanding)
        \draw[up] (b.east) -| node[pos=0.25, above, scale=0.8, text=uclblue] {Up} (d2.south);
        \draw[up] (d2.east) -| node[pos=0.25, above, scale=0.8, text=uclblue] {Up} (d1.south);
        \draw[up] (d1.east) -| node[pos=0.25, above, scale=0.8, text=uclblue] {Up} (d0.south);
        
        % Skip Connections
        \draw[skip] (e0) -- node[above, scale=0.8, text=darkgreen] {Skip + Concat} (d0);
        \draw[skip] (e1) -- node[above, scale=0.8, text=darkgreen] {Skip + Concat} (d1);
        \draw[skip] (e2) -- node[above, scale=0.8, text=darkgreen] {Skip + Concat} (d2);
        
    \end{tikzpicture}
    }
    \end{center}
\end{frame}

%================================================================%
\section{Direct Supervised Learning (End-to-End) \& Its Pitfalls}
%================================================================%

\begin{frame}{The Temptation of End-to-End Learning}
    \only<1-2>{Neural networks excel at learning complex input--output mappings directly from data. It is tempting to apply this to \textbf{any} process, including inverse problems:
    \begin{itemize}
        \item Image denoising, superresolution, deblurring, phase retrieval, CT reconstruction, MRI reconstruction, \ldots
    \end{itemize}}
    
    \onslide<2->
    The idea is simple: given a forward model
    \begin{align*}
        Ax = y,
    \end{align*}
    where $A \in \mathbb{R}^{m \times n}$ is a known (forward) operator, $x \in \mathbb{R}^n$ is the object to recover, and $y \in \mathbb{R}^m$ is the observed noisy data, why not simply learn an approximation to the Moore--Penrose pseudo-inverse $A^\dagger$?
    
    \only<3->{\begin{block}{Moore--Penrose Pseudo-Inverse \protect\citep{ehn}}
        For any matrix $A$, the Moore--Penrose inverse is defined as $$ A^\dagger y := \argmin_x \left\{ \| x \| \, \left| \,  x \in \argmin_{\tilde{x}} \|A\tilde{x} - y\|^2 \right. \right\} = \argmin_x \left\{ \| x \| \, \left| \,  A^\top A x = A^\top y \right. \right\} \, . $$ 
        If $A$ has full column rank, $A^\dagger = (A^\top A)^{-1} A^\top$.
    \end{block}}
\end{frame}

\begin{frame}{End-to-End Learning for Inversion}
    \textbf{Idea:} Train a neural network $f_w$ to map noisy observations directly to recovered objects \citep{arridge2019solving}:
    \begin{align*}
        \min_w \frac{1}{N} \sum_{i=1}^N \left\| f_w(y_i^\delta) - x_i^\dagger \right\|^2,
    \end{align*}

    \vspace{-0.15cm}
    where $(y_i^\delta, x_i^\dagger)$ are training pairs (noisy data and ground truth solutions), and we hope $f_w \approx A^\dagger$.\pause
    
    \textbf{Attractions:}
    \begin{enumerate}
        \item \textbf{Simplicity:} No inverse problem theory required; just standard supervised learning.
        \item \textbf{Speed:} Single forward pass at inference; no iterative solves needed.
        \item \textbf{Empirical Performance:} Often achieves high PSNR/SSIM scores on restricted test sets matching the training distribution.
    \end{enumerate}
    
    \pause
    \begin{center}
        \textbf{However, direct end-to-end inversion is fundamentally fragile and can fail catastrophically.}
    \end{center}
\end{frame}

\begin{frame}{Pitfall 1: Lack of Data Consistency}
    The network output $\hat x = f_w(y^\delta)$ is a prediction learned from noisy training data. It is \textbf{not constrained} to satisfy the forward model.
    
    \begin{alertblock}{Forward Model is Ignored}
        There is \textbf{no guarantee} that $A \hat x_i = y_i$, where $y_i = A x^\dagger_i$.
    \end{alertblock}
    
    \begin{itemize}
        \item The network may generate a visually plausible image that violates the forward model, i.e. $\|A \hat x - y^\delta\|^2$ can be arbitrarily large.
        \item \textbf{Consequence:} In medical imaging, the network might remove a small tumour (treating it as noise) or hallucinate healthy tissue where abnormalities exist, since the training loss never enforces consistency with $y^\delta$.
    \end{itemize}
\end{frame}

\begin{frame}{Pitfall 2: The Lipschitz Dilemma \& Adversarial Vulnerability}
    \textbf{Mathematical Reality:} For matrices $A$ with fast decaying singular values, $A^\dagger$ is \emph{ill-conditioned}; small changes in $y$ can cause arbitrarily large changes in $A^\dagger y$.
    
    \textbf{Network Limitation:} Neural networks are continuous functions with a finite Lipschitz constant $L$, i.e.,
    $$\|f_w(y^\delta + \eta) - f_w(y^\delta)\| \leq L \|\eta\| \, . $$
    
    \textbf{The Conflict:} To approximate a highly ill-conditioned operator $A^\dagger$ on the training set, the network is forced to learn an astronomically large Lipschitz constant $L$. While it may perform well on the training data manifold, this massive $L$ creates vulnerabilities off the manifold.
    
    \pause
    \vspace{0.5em}
    \begin{alertblock}{Adversarial Instability \protect\citep{antun2020instabilities}}
        Because $L$ is so large, a worst-case perturbation $\eta$ with tiny $\|\eta\|$ can cause $\|f_w(y^\delta + \eta) - f_w(y^\delta)\|$ to be enormous, producing severe reconstruction artefacts.
    \end{alertblock}
\end{frame}

\begin{frame}{Pitfall 3: The "Opaque" Implicit Bias}
    \textbf{Classical Mathematical Methods:} The assumptions about the solution are explicit. We define exactly what properties we expect a recovered object to have (e.g., piece-wise smoothness or sharp edges).
    
    \textbf{End-to-End Learning:} The assumptions are implicit and hidden. They are determined solely by the statistics of the training data.
    
    \begin{itemize}
        \item The network learns to fill in the null space of $A$ (the space where $Av = 0$) using whatever patterns it observed in training. This is \emph{memorisation}.
        \item \textbf{Generalisation Failure:} When a test case contains a rare or unseen feature (e.g., an atypical tumour morphology), the network has never learned a representation for it. Instead, it hallucinates a "typical" feature based on its training bias, replacing the truth.
    \end{itemize}
\end{frame}

\begin{frame}{Summary of Dangers: Classical Methods vs. Deep Learning}
    \begin{table}
        \centering
        \begin{tabular}{l|c|c}
             & \textbf{Classical Methods} & \textbf{Deep Learning (End-to-End)} \\
             \hline
             \textbf{Data Consistency} & Guaranteed ($Ax \approx y$) & Not Guaranteed \\
             \textbf{Stability} & Mathematically bounded & Unstable (High Lipschitz needed) \\
             \textbf{Assumptions} & Explicit (Hand-crafted) & Opaque (Training data distribution) \\
             \textbf{Image Quality} & Limited (Oversimplified) & High (Realistic textures) \\
        \end{tabular}
    \end{table}
    
    \vspace{1em}
    \textbf{Goal for the Course:} Learn to combine the reliability of mathematical forward models with the expressivity of deep learning. We will explore topics such as unrolled networks, implicit deep learning, and advanced generative models.
\end{frame}

\begin{frame}{Summary \& Preparation for the Coding Challenge}
    \textbf{Review:} From perceptrons to the pitfalls of black-box inversion, ending with continuous and implicit models.
    
    \vspace{1em}
    \textbf{Coding Challenge:}
    \begin{itemize}
        \item Recovering sharp images from degraded measurements using PyTorch/DeepInverse.
        \item Baseline: Simple mathematical inversion.
        \item Constraints: Max 100 MB model size, CPU inference $< 20$ mins.
    \end{itemize}
\end{frame}

\begin{frame}[allowframebreaks]{References}
    \bibliographystyle{plainnat}
    \bibliography{references}
\end{frame}

\end{document}
